\documentclass[12pt]{article}

\usepackage[utf8]{inputenc}
\usepackage{amsmath, amssymb, amsthm}
\usepackage{geometry}
\geometry{a4paper, margin=2.5cm}

\title{Teorema de Unicitate a Probabilităților}
\author{}
\date{}

\theoremstyle{plain}
\newtheorem{theorem}{Teoremă}

\begin{document}

\maketitle

\begin{theorem}[Teorema de unicitate a probabilităților]
Fie $(\Omega, \mathcal{F})$ un spațiu măsurabil și fie 
$\mathbb{P}_1, \mathbb{P}_2 : \mathcal{F} \to [0,1]$ două probabilități.
Fie $\mathcal{C} \subseteq \mathcal{F}$ o familie de mulțimi.

\begin{enumerate}
    \item Dacă $\mathbb{P}_1$ și $\mathbb{P}_2$ coincid pe $\mathcal{C}$, atunci ele coincid pe lambda-sistemul generat de $\mathcal{C}$:
    

\[
        \mathbb{P}_1(A) = \mathbb{P}_2(A)
        \quad \text{pentru toate } A \in \lambda(\mathcal{C}).
    \]



    \item Dacă, în plus, $\mathcal{C}$ este un $\pi$-sistem, atunci
    

\[
        \lambda(\mathcal{C}) = \sigma(\mathcal{C}),
    \]


    iar probabilitățile coincid pe întreaga sigma-algebră generată:
    

\[
        \mathbb{P}_1(A) = \mathbb{P}_2(A)
        \quad \text{pentru toate } A \in \sigma(\mathcal{C}).
    \]


\end{enumerate}
\end{theorem}

\end{document}
